\booksection{Κεφάλαιο 1 Παράρτη$\mu$α: Ειδικό θέ$\mu$α λειτουργικού συστή$\mu$ατος HarmonyOS}

\chapterintro{Το HarmonyOS είναι λειτουργικό σύστη$\mu$α αυτοανεπτυγ$\mu$ένο από την Huawei, από το 2019 έχει περάσει πολλές $\mu$εγάλες εκδόσεις. Δενείναι «ακό$\mu$α $\mu$ια αντικατάσταση του Android», αλλά ένα νέο είδος $\mu$ε $\mu$οναδική φιλοσοφία σχεδιασ$\mu$ού. Ας το γνωρίσου$\mu$ε σε βάθος.}

\booksubsection{I. Τιείναι το HarmonyOS; Γιατί είναι διαφορετικό;}

Η βασική φιλοσοφία του HarmonyOS είναι «σύνδεση όλων». Σε αντίθεση $\mu$ε τα παραδοσιακά λειτουργικά που εξυπηρετούν κυρίως $\mu$ία συσκευή, το HarmonyOS από τον σχεδιασ$\mu$ό τουλα$\mu$βάνει υπόψη συνεργασία πολλαπλών συσκευών.

\textbf{Κατανε$\mu$η$\mu$ένη αρχιτεκτονική}είναι η ψυχή του HarmonyOS. Απλά, $\mu$πορεί ναεικονικό κάνει πολλές συσκευές σε ένα «σούπερ τερ$\mu$ατικό». Το κινητό, tablet, υπολογιστής, ρολόι, TV σας δενείναι πλέον $\mu$ε$\mu$ονω$\mu$ένες οντότητες, αλλά ένα ολόκληρο «όργανο», που$\mu$πορούν α$\mu$οιβαία καλεί τις ικανότητες του άλλου.

Παράδειγ$\mu$α: βλέπετε βίντεο στο κινητό και θέλετε στην $\mu$εγάλη οθόνη. Παραδοσιακά είναι προβολή, στο HarmonyOS $\mu$πορείτε να $\mu$εταφέρετε το βίντεο «ροή» στην TV, το κινητό γίνεται τηλεχειριστήριο. Δενείναι απλόπροβολή, αλλά η εφαρ$\mu$ογή τρέχει διασυσκευαστικό.

\booksubsection{II. HarmonyOS 4.3: ώρι$\mu$η έκδοση βασισ$\mu$ένη στο AOSP}

Το HarmonyOS 4.3 είναι η τρέχουσα έκδοσηευρέως αναπτυχγ$\mu$ένο στα κινητά Huawei. Ση$\mu$ειώστε ότι το HarmonyOS 4.3 συ$\mu$βατό το Android Open Source Project (AOSP) framework, γι' αυτό$\mu$πορεί πλήρως συ$\mu$βατό $\mu$ε Android εφαρ$\mu$ογές — οι περισσότερες App που κατεβάζετε από την αγορά εφαρ$\mu$ογών Huawei είναι ουσιαστικά Android apps, που τρέχουν $\mu$έσω του compatibility layer του HarmonyOS. Ό$\mu$ως το HarmonyOS 4.3 δενείναι απλάεπικάλυψη Android, στο AOSP προσθέτει $\mu$εγάλο όγκο αυτοανεπτυγ$\mu$ένο κατανε$\mu$η$\mu$ένες δυνατότητες, Ark Engine και άλλα τεχνολογίες.

\subsubsection*{Βασικά χαρακτηριστικά:}

\begin{itemize}
  \item \textbf{Συ$\mu$βατότητα $\mu$ε Android εφαρ$\mu$ογές}: Επειδή βασίζεται στο AOSP, το HarmonyOS 4.3 $\mu$πορεί να τρέξει σχεδόν όλες τις εφαρ$\mu$ογές Android, οι χρήστες σχεδόν δεν παρατηρούν διαφορά.
  \item \textbf{Ark Engine}: ο αυτοανεπτυγ$\mu$ένος $\mu$εταγλωττιστής και χρόνος εκτέλεσης της Huawei, κάνει τις εφαρ$\mu$ογές να τρέχουν πιο ο$\mu$αλά και $\mu$ε λιγότερη κατανάλωση. Σε σύγκριση $\mu$ε προηγού$\mu$ενες εκδόσεις, η ο$\mu$αλότητα έχει βελτιωθεί ση$\mu$αντικά.
  \item \textbf{Υπερ-τερ$\mu$ατικό}: Στο Κέντρο Ελέγχου, $\mu$πορείτε να δείτε διαισθητικά τις άλλες κοντινές συσκευές HarmonyOS και να τις συνδέσετε σύροντάς τες. Το τηλέφωνο $\mu$πορεί να χρησι$\mu$οποιήσει το πληκτρολόγιο και το ποντίκι του υπολογιστή, και το tablet την κά$\mu$ερα του τηλεφώνου — όλα πολύ εύκολα.
  \item \textbf{Κάρτα υπηρεσίας}: Παρό$\mu$οιο $\mu$ε τα widgets του iOS, αλλά πιο ευέλικτο. Σύρετε προς τα πάνω το εικονίδιο της εφαρ$\mu$ογής για να αναπτυχθεί $\mu$ια κάρτα $\mu$ε βασικές πληροφορίες, χωρίς να ανοίξετε την εφαρ$\mu$ογή.
  \item \textbf{Λειτουργία καθαρής εγκατάστασης}: επιτρέπει $\mu$όνο εγκατάσταση εφαρ$\mu$ογών από το επίση$\mu$ο κατάστη$\mu$α εφαρ$\mu$ογών, $\mu$ειώνοντας τον κίνδυνο κακόβουλου λογισ$\mu$ικού.
  \item \textbf{Κέντρο απορρήτου}: κεντρική διαχείριση δι και ω$\mu$άτων εφαρ$\mu$ογών, προβολή αρχείου δραστηριότητας.
\end{itemize}

\subsubsection*{Χαρακτηριστικά διεπαφής χρήστη:}

\begin{itemize}
  \item \textbf{Κλείδω$\mu$α οθόνης περιοδικού}: η οθόνη κλειδώ$\mu$ατος ε$\mu$φανίζει ό$\mu$ορφες εικόνες και άρθρα, σύρετε για να $\mu$πείτε στην επιφάνεια εργασίας.
  \item \textbf{Παγκόσ$\mu$ια κάρτα}: τα εικονίδια επιφάνειας εργασίας $\mu$πορούν να επεκταθούν σε κάρτες, ε$\mu$φανίζοντας περισσότερες πληροφορίες.
  \item \textbf{Σούπερ κέντρο $\mu$εταφόρτωσης}: $\mu$ια προσωρινή περιοχή αποθήκευσης, $\mu$πορείτε να αποθηκεύσετε προσωρινά κεί$\mu$ενο, εικόνες, αρχεία και $\mu$ετά να τα σύρετε σε οποιαδήποτε εφαρ$\mu$ογή. Η δια$\mu$οιρασ$\mu$ός $\mu$εταξύ εφαρ$\mu$ογών γίνεται απίστευτα ο$\mu$αλή.
  \item \textbf{Έξυπνη αναζήτηση}: καθολική αναζήτηση, $\mu$πορεί να αναζητήσει εφαρ$\mu$ογές, αρχεία, ρυθ$\mu$ίσεις, ιστοσελίδες.
\end{itemize}

\booksubsection{III. HarmonyOS NEXT (HarmonyOS 6): η αληθινή αφετηρία του αυτοανεπτυγ$\mu$ένου λειτουργικού της Huawei}

Το HarmonyOS NEXT (HarmonyOS NEXT 5.0) έχει θε$\mu$ελιώδη διαφορά από το HarmonyOS 4.3 — υιοθετεί τον αυτοανεπτυγ$\mu$ένο $\mu$ικροπυρήνα Harmony της Huawei, δεν βασίζεται πλέον στον πυρήνα Linux και τον κώδικα AOSP, και\textbf{Δεν υποστηρίζει πλέον ά$\mu$εσα εφαρ$\mu$ογές Android}. Αυτό είναι το κρίσι$\mu$ο βή$\mu$α του Harmony από “$\mu$εταβατική $\mu$ορφή συ$\mu$βατή $\mu$ε Android” προς “πλήρως ανεξάρτητο λειτουργικό σύστη$\mu$α”.

\subsubsection*{Ση$\mu$αντικές αλλαγές:}

\begin{itemize}
  \item \textbf{Πλήρως αυτοανεπτυγ$\mu$ένος πυρήνας}: το HarmonyOS NEXT χρησι$\mu$οποιεί τον αυτοανεπτυγ$\mu$ένο πυρήνα της Huawei, δεν περιέχει πλέον καθόλου κώδικα AOSP. Το σύστη$\mu$α είναι πιο λιτό, πιο ασφαλές και πιο αποδοτικό.
  \item \textbf{Δεν υποστηρίζει πλέον Android εφαρ$\mu$ογές}: αυτή είναι η $\mu$εγαλύτερη αλλαγή. Το HarmonyOS NEXT $\mu$πορεί να τρέξει $\mu$όνο εγγενείς εφαρ$\mu$ογές HarmonyOS, δεν $\mu$πορεί να εγκαταστήσει ή να τρέξει αρχεία APK. Πριν αναβαθ$\mu$ίσετε, επιβεβαιώστε ότι οι App που χρησι$\mu$οποιείτε έχουν ήδη εγγενείς εκδόσεις HarmonyOS.
  \item \textbf{Οικοσύστη$\mu$α εγγενές εφαρ$\mu$ογών}: η Huawei προωθεί ενεργά τους developers στη $\mu$εταφορά εφαρ$\mu$ογών στην εγγενή πλατφόρ$\mu$α HarmonyOS. Οι κύριες App (WeChat, Alipay, Taobao) έχουν ήδη εγγενείς εκδόσεις HarmonyOS, αλλά κάποιες niche εφαρ$\mu$ογές $\mu$πορεί να $\mu$ην είναι ακό$\mu$α προσαρ$\mu$οσ$\mu$ένες.
  \item \textbf{Star Shield Ασφαλής αρχιτεκτονική}: ολοκαίνουργιο σύστη$\mu$α ασφαλείας, προστασία end-to-end από hardware σε software. Υποστηρίζει πιο λεπτο$\mu$ερή έλεγχο δι και ω$\mu$άτων, π.χ. «θολή τοποθεσία» — η εφαρ$\mu$ογή βλέπει $\mu$όνο κατά προσέγγιση περιοχή, όχι ακριβές ση$\mu$είο.
  \item \textbf{Βαθιά ενσω$\mu$άτωση AI}: Δυνατότητες AI σε επίπεδο συστή$\mu$ατος, όχι $\mu$όνο σε $\mu$ε$\mu$ονω$\mu$ένες εφαρ$\mu$ογές. Για παράδειγ$\mu$α, οι AI υπότιτλοι $\mu$εταφράζουν βίντεο σε πραγ$\mu$ατικό χρόνο, το AI αφαίρεσης φόντου αναγνωρίζει έξυπνα το θέ$\mu$α.
  \item \textbf{Διασυσκευή απρόσκοπτη ροή}: Πιο προηγ$\mu$ένη κατανε$\mu$η$\mu$ένη ικανότητα. Η κατάσταση των εφαρ$\mu$ογών $\mu$πορεί να $\mu$εταβαίνει απρόσκοπτα $\mu$εταξύ συσκευών — το έγγραφο που επεξεργάζεστε στο τηλέφωνο ανοίγει αυτό$\mu$ατα στον υπολογιστή.
  \item \textbf{Atomic Service}: Ελαφριές υπηρεσίες που δεν χρειάζονται εγκατάσταση. Για παράδειγ$\mu$α, σαρώστε ένα QR code και ανοίγει απευθείας η λειτουργία χωρίς λήψη εφαρ$\mu$ογής.
\end{itemize}

\begin{tcbwarn}{Πριν αναβαθ$\mu$ίσετε στο HarmonyOS NEXT —απαραίτητο διάβασ$\mu$α}
Αφότου αναβαθ$\mu$ίσετε στο HarmonyOS NEXT, το κινητό σας δεν $\mu$πορεί να τρέξει Android εφαρ$\mu$ογές. Πριν την αναβάθ$\mu$ιση πρέπει οπωσδήποτε επιβεβαιώστε: (1) οι App πουσυχνά χρησι$\mu$οποιούνται εγγενές HarmonyOS έκδοση· (2) στηεργασία ή ζωή σας δεν υπάρχουν σενάρια πουπρέπει εξαρτάται Android εφαρ$\mu$ογές. Μετά την αναβάθ$\mu$ιση δεν $\mu$πορεί να επιστρέψει στο HarmonyOS 4.3.

\end{tcbwarn}

\subsubsection*{Αλλαγές λογικής λειτουργίας:}

\begin{itemize}
  \item \textbf{Πλοήγηση $\mu$ε χειρονο$\mu$ίες}: σύρετε από τις αριστερές/δεξιές άκρες για επιστροφή, σύρετε από κάτω προς τα πάνω για την επιφάνεια εργασίας, σύρετε προς τα πάνω και κρατήστε για πολλαπλές εργασίες.
  \item \textbf{Έξυπνα πολλαπλά παράθυρα}: υποστηρίζει split screen, floating window, parallel view και άλλες $\mu$ορφές πολλαπλών εργασιών.
  \item \textbf{Changlian κλήση}: Η υπηρεσία βιντεοκλήσεων της Huawei, $\mu$ε κοινή χρήση οθόνης, φίλτρα ο$\mu$ορφιάς και θόλωση φόντου.
\end{itemize}

\booksubsection{IV. Οικοσύστη$\mu$α HarmonyOS: συσκευές και εφαρ$\mu$ογές}

Το HarmonyOS δενείναι $\mu$όνο σύστη$\mu$α κινητού, αλλά οικοσύστη$\mu$α που καλύπτει πολλές συσκευές.

\subsubsection*{Υποστηριζό$\mu$ενοι τύποι συσκευών:}

\begin{itemize}
  \item \textbf{Κινητά και tablets}: Η κύρια συσκευή, οι πιο πλήρεις λειτουργίες.
  \item \textbf{Smartwatches και fitness bands}: λειτουργεί ανεξάρτητα ή σε συνεργασία $\mu$ε το κινητό.
  \item \textbf{Smart οθόνη (TV)}: κέντρο οικιακής ψυχαγωγίας.
  \item \textbf{Έξυπνο ηχείο}: πύλη φωνητικής αλληλεπίδρασης.
  \item \textbf{Αυτοκίνητο}: HarmonyOS Cockpit — ε$\mu$πειρία έξυπνου cockpit.
  \item \textbf{Έξυπνο σπίτι}: ψυγείο, κλι$\mu$ατιστικό, πλυντήριο κ.λπ.
\end{itemize}

\subsubsection*{Τρέχουσα κατάσταση οικοσυστή$\mu$ατος εφαρ$\mu$ογών:}

Το App Gallery του HarmonyOS έχει εκατο$\mu$$\mu$ύρια εφαρ$\mu$ογές, $\mu$ε τις κυριότερες App να είναι διαθέσι$\mu$ες. Η Huawei προωθεί ενεργά τους προγρα$\mu$$\mu$ατιστές να δη$\mu$ιουργούν εγγενείς εφαρ$\mu$ογές HarmonyOS. Μέχρι τα $\mu$έσα του 2026, WeChat, Alipay, Taobao, Douyin, Meituan, JD.com και άλλες κορυφαίες εφαρ$\mu$ογές έχουν κυκλοφορήσει εγγενείς εκδόσεις HarmonyOS, και η καθη$\mu$ερινή χρήση είναι σχεδόν χωρίς ε$\mu$πόδια. Ορισ$\mu$ένες $\mu$ικρές εφαρ$\mu$ογές και επαγγελ$\mu$ατικά εργαλεία ίσως προσαρ$\mu$ογή, επο$\mu$ένως πριν την αναβάθ$\mu$ιση Συστήνεται να επιβεβαιώσετε ότι οι App που χρησι$\mu$οποιείτε συχνά υποστηρίζονται.

\booksubsection{V. Πρακτικές συ$\mu$βουλές HarmonyOS}

\begin{itemize}
  \item \textbf{Υπερ-τερ$\mu$ατικό}: Σύρετε προς τα κάτω από την πάνω δεξιά γωνία για να ανοίξετε το Κέντρο Ελέγχου, πατήστε "Super Device" — ε$\mu$φανίζονται οι κοντινές συσκευές, σύρετε για σύνδεση.
  \item \textbf{Σούπερ κέντρο $\mu$εταφόρτωσης}: παρατετα$\mu$ένο πάτη$\mu$α σε κεί$\mu$ενο/εικόνα/αρχείο, επιλέξτε "Αντιγραφή στον ενδιά$\mu$εσο σταθ$\mu$ό", $\mu$ετά σε άλλη εφαρ$\mu$ογή παρατετα$\mu$ένο πάτη$\mu$α για επικόλληση ή σύρετε έξω.
  \item \textbf{Κάρτα υπηρεσίας}: Σύρετε προς τα πάνω το εικονίδιο της εφαρ$\mu$ογής για να ε$\mu$φανιστεί η κάρτα, πατήστε παρατετα$\mu$ένα για αλλαγή $\mu$εγέθους ή προσθήκη στην αρχική οθόνη.
  \item \textbf{Έξυπνα πολλαπλά παράθυρα}: σύρετε από την αριστερή ή δεξιά άκρη της οθόνης προς τα $\mu$έσα και κρατήστε, ανοίξτε την πλαϊνή $\mu$πάρα, πατήστε την εφαρ$\mu$ογή για να ανοίξει σε αιωρού$\mu$ενο παράθυρο.
  \item \textbf{Προστασία ιδιωτικότητας}: Ρυθ$\mu$ίσεις > Απόρρητο, $\mu$πορείτε να δείτε το αρχείο χρήσης δι και ω$\mu$άτων εφαρ$\mu$ογών, να ορίσετε ασαφή τοποθεσία κ.λπ.
\end{itemize}

\begin{tcbwarn}{Ση$\mu$ειώσεις}
Το HarmonyOS ενη$\mu$ερώνεται και εξελίσσεται συνεχώς, διαφορετικές εκδόσεις $\mu$πορεί να έχουν$\mu$ερικά διαφορές. Η περιγραφή εδώ βασίζεται στις κύριες λειτουργίες του HarmonyOS 4.3 και 6, οι συγκεκρι$\mu$ένες λειτουργίες ανατρέξτε την πραγ$\mu$ατική έκδοση του κινητού σας. Προτείνεται να διατηρείτε το σύστη$\mu$α ενη$\mu$ερω$\mu$ένο, για να έχετε τις νεότερες λειτουργίες και ταασφάλεια patches.

\end{tcbwarn}

\begin{tcbpractice}{Δοκι$\mu$άστε τώρα: Βιώστε τα χαρακτηριστικά του HarmonyOS}
\begin{enumerate}
  \item Ανοίξτε το κέντρο ελέγχου, βρείτε το «Σούπερ τερ$\mu$ατικό», δείτε αν υπάρχουν άλλες HarmonyOS συσκευές γύρω.
  \item Δοκι$\mu$άστε το «Σούπερ κέντρο $\mu$εταφόρτωσης»: αντιγράψτε κεί$\mu$ενο στο κέντρο $\mu$εταφόρτωσης, $\mu$ετά σύρετέ τοέξω στις ση$\mu$ειώσεις.
  \item Σύρετε προς τα πάνω $\mu$ερικάσυχνά χρησι$\mu$οποιού$\mu$ενα εικονίδια App, δείτε τι ε$\mu$φανίζουν οικάρτα υπηρεσίας.
  \item Πηγαίνετε Ρυθ$\mu$ίσεις > Απόρρητο, δείτεπρόσφατα ποιες εφαρ$\mu$ογέςχρήση ευαίσθητων δι και ω$\mu$άτων.
\end{enumerate}

\end{tcbpractice}
